En los últimos años, la integración de herramientas basadas en Inteligencia Artificial Generativa (IAG) ha empezado a cambiar flujos de trabajos en el desarrollo de videojuegos, abarcando desde la generación de assets visuales y elementos narrativos hasta la asistencia en la programación del juego. La literatura encontrada destaca el potencial de la IAG para optimizar procesos creativos, reducir tiempos de producción y habilitar nuevas dinámicas de interacción entre el jugador y el sistema, su implementación ha creado un conflicto dentro de la industria. Por un lado, la investigación técnica destaca la eficiencia y el diseño procedimental asistido por IA; por el otro lado, la recepción por parte de los jugadores revela rechazo a este tipo de tecnologías, creando escepticismo ético y la percepción de pérdida de autenticidad debido al contenido generado por estos modelos. Esta brecha entre la capacidad técnica del desarrollo y la percepción subjetiva del usuario plantea la necesidad de examinar de manera rigurosa cómo la presencia de la IA impacta directamente en la experiencia del jugador.  

Esta resistencia por parte del público ha dejado de ser una especulación y se ha hecho visible en la industria actual, observable en el caso del videojuego \textit{Clair Obscure: Expedition 33}, donde la revelación o sospecha del uso de IA generativa en la producción derivó en un rechazo inmediato por parte de la comunidad de jugadores y en un debate sobre la legitimidad del contenido \cite{hollyhead2026clair}. Este tipo de reacciones se alinea con lo documentado en la literatura sobre percepción pública de la IA; investigaciones sobre la evaluación de contenido sintético demuestran que los usuarios tienden a penalizar severamente el contenido cuando son conscientes de su origen automatizado, afectando directamente la credibilidad y el valor percibido del producto en comparación con las creaciones humanas \cite{huschens2023}. Sin embargo, este sesgo negativo no es homogéneo dentro de la comunidad. En el estudio de Quantic Foundry \cite{yee2025genai}, la postura de los jugadores hacia esta herramienta en videojuegos oscila considerablemente según factores demográficos y según las motivaciones del jugador, encontrando una mayor aversión entre los usuarios con un alto compromiso creativo, frente a sectores más permisivos centrados únicamente en mecánicas o rendimiento técnico.

Por otro lado, desde la perspectiva del desarrollo de software y videojuegos, se evidencia un interés en el uso de la IA para mantener un nivel de eficiencia con la generación de contenido procedimental \cite{filipovic2023}. Herramientas como las basadas en aprendizaje por refuerzo no solo han permitido evaluar y optimizar mecánicas de juego complejas en etapas tempranas de diseño \cite{gonzalez2023mechanic}, sino también asistir en tareas pesadas de programación y generación de recursos \cite{li2024}. No obstante, los propios creadores se encuentran con el dilema que mientras la IAG acelera la ideación y el prototipado, emergen serias inquietudes éticas y profesionales vinculadas a la pérdida de originalidad, la dependencia tecnológica y la posible devaluación de la autoría humana \cite{alharthi2025}. Este conflicto en el desarrollo se traslada al consumidor final; haciendo que la aceptación de la IA se deba a la función específica que esta cumpla \cite{hsu2026ai}. Mientras que los jugadores son más permisivos hacia aplicaciones invisibles o de soporte, manifiestan una marcada resistencia cuando la IA sustituye procesos creativos visibles como el diseño de arte, guión o dirección estética general del juego.

Al examinar cómo estas percepciones y usos técnicos se traducen en la vivencia interactiva del usuario, la literatura sobre la experiencia de usuario (UX) revela que la percepción de la presencia de IA no es un factor neutral, sino una variable que altera la arquitectura de la experiencia de juego. Tradicionalmente, la inmersión y la respuesta cognitiva del usuario han dependido del \textit{game feel} y de la percepción de intencionalidad en el diseño interactivo \cite{swink2008gamefeel}. Se ha podido ver que en revisiones hechas como la de Lao \cite{lao2025} y estudios sobre contenido generado \cite{bazzaz2026imitation} demuestran la creencia de que un elemento fue producido por un algoritmo en lugar de un ser humano altera negativamente dimensiones psicológicas clave, como la inmersión, el valor percibido del reto y la conexión emocional con el mundo del juego. 

A partir de toda la literatura investigada, se evidencia un amplió análisis centrado en el uso de IA para el desarrollo de videojuegos. No obstante, existen vacíos teóricos en este campo; en particular, la respuesta actitudinal del consumidor final frente al uso de esta tecnología en videojuegos no ha sido estudiada con la misma profundidad. Si bien investigaciones como la de Zhang y Gosline \cite{zhang2023aicreated} han evaluado la percepción pública ante contenidos de diverso origen (humano versus AI) en ámbitos como la publicidad y los productos comerciales, existe una escasez de evidencia focalizada en la industria del videojuego. En consecuencia, la presente investigación aborda este problema analizando la percepción de la comunidad de jugadores sobre la utilización de IAG en el desarrollo de videojuegos, con el propósito de identificar los factores psicológicos, éticos y de experiencia de usuario que subyacen al rechazo manifestado hacia el contenido generado por IA en esta industria.